%!TEX root = ../main.tex
\chapter{QMC-Based Deep Learning Algorithm}
\label{chapter6}

The goal of this chapter is to mathematically analyze the function approximation
problem and find bounds on the generalization gap of \ac{dnn} surrogate models
obtained by low-discrepancy training data. The analysis is based on the
Hardy--Krause variation and the Koksma-Hlawka inequality introduced in
Chapter~\ref{chapter3}.

\section{Preliminaries}

For the sake of completeness, we start with some results on approximations of
variations defined in Section~\ref{sec:variation}. First we recall an important
identity for the alternating sum.

The following results come from \cite[Section 9]{owen2005multidimensional}. For
completeness, the conceptual proofs are given here in full.

\begin{prop} \ \\
\label{prop:variation-approximation}
    Let $f$ be a function defined on the hypercube $[a,b]\subset \mathbb{R}^s$
    with existing and integrable derivative $\frac{\partial^s}{\partial x_1
    \cdots \partial x_s} f$. Then, the following holds for the alternating sum:
    \begin{equation}
        \int\limits_{[a,b]} \frac{\partial^s}{\partial x_1 \cdots \partial x_s} 
        f(x) \, dx = \Delta(f;a,b)
    \end{equation}
    where $\Delta(f;a,b)$ is the alternating sum from
    Definition~\ref{def:alternating_sum}.
\end{prop}

\begin{proof} \ \\
The proof follows by induction on the dimension $s$.

\textbf{Base case ($s=1$).}
\begin{equation*}
    \int\limits_{a}^{b} \frac{\partial f(x)}{\partial x} \, dx
    = f(b) - f(a) = \Delta(f;a,b).
\end{equation*}
Thus, the statement holds for $s=1$.

\textbf{Induction hypothesis.}\\
Assume the statement holds for dimension $s-1$.

\textbf{Induction step ($s-1 \to s$).}\\
Write $x = (x',x_s)$ with $x' \in [a',b'] := \prod_{j=1}^{s-1} [a_j,b_j]$.
First, apply the one-dimensional fundamental theorem of calculus in the last
coordinate:
\begin{equation*}
\int\limits_{a_s}^{b_s} 
\frac{\partial^s f(x',t)}{\partial x_1 \cdots \partial x_s} \, dt
= \frac{\partial^{s-1} f(x',b_s)}{\partial x_1 \cdots \partial x_{s-1}}
 - \frac{\partial^{s-1} f(x',a_s)}{\partial x_1 \cdots \partial x_{s-1}}.
\end{equation*}
Now integrate over $x'$ and use Fubini’s theorem:
\begin{equation*}
\int\limits_{[a,b]} 
\frac{\partial^s f(x)}{\partial x_1 \cdots \partial x_s} \, dx
= \int\limits_{[a',b']} \frac{\partial^{s-1} f(x',b_s)}{\partial x_1 \cdots \partial x_{s-1}} \, dx'
 - \int\limits_{[a',b']} \frac{\partial^{s-1} f(x',a_s)}{\partial x_1 \cdots \partial x_{s-1}} \, dx'.
\end{equation*}
We now apply the induction hypothesis to both integrals (for the functions
$g_1(x') := f(x',b_s)$ and $g_2(x') := f(x',a_s)$):
\begin{equation*}
\int\limits_{[a',b']} \frac{\partial^{s-1} f(x',b_s)}{\partial x_1 \cdots \partial x_{s-1}} \, dx'
= \Delta\big(f(\cdot,b_s);a',b'\big),
\end{equation*}
\begin{equation*}
\int\limits_{[a',b']} \frac{\partial^{s-1} f(x',a_s)}{\partial x_1 \cdots \partial x_{s-1}} \, dx'
= \Delta\big(f(\cdot,a_s);a',b'\big).
\end{equation*}
Thus,
\begin{equation*}
\int\limits_{[a,b]} 
\frac{\partial^s f(x)}{\partial x_1 \cdots \partial x_s} \, dx
= \Delta\big(f(\cdot,b_s);a',b'\big)
 - \Delta\big(f(\cdot,a_s);a',b'\big).
\end{equation*}
This is exactly the recursive definition of the corner difference in $s$
dimensions. Hence, the statement holds for $s$.
\end{proof}

This yields another result for the variation in the sense of Vitali on functions
with existing and integrable derivative.

\begin{cor} \ \\
\label{cor:variation-approximation}
    Let $f$ be a function defined on the hypercube $[a,b]\subset \mathbb{R}^s$ with existing and
    integrable derivative $\frac{\partial^s}{\partial x_1 \cdots \partial x_s}
    f$. Then, the following holds for the variation in the sense of Vitali:
    \begin{equation*}
        V_{[a,b]}(f) \leq \int\limits_{[a,b]} \left| \frac{\partial^s}{\partial x_1 \cdots \partial x_s} f(x) \right| \, dx.
    \end{equation*}
\end{cor}

\begin{proof} \ \\
Recalling the definition of variation with respect to a ladder from
Definition~\ref{def:variation_ladder} and applying
Proposition~\ref{prop:variation-approximation} yields

\begin{equation*}
    V_{\mathcal{Y}}(f) =  \sum_{y \in \mathcal{Y}} \big| \Delta(f; y, y_+) \big| = \sum_{y \in \mathcal{Y}} \bigg| \int\limits_{[y,y_+]}\frac{\partial^s}{\partial x_1 \cdots \partial x_s} f(x) \, dx \bigg|.
\end{equation*}

Applying the triangle inequality then leads to

\begin{equation*}
    \sum_{y \in \mathcal{Y}} \bigg| \int\limits_{[y,y_+]}\frac{\partial^s}{\partial x_1 \cdots \partial x_s} f(x) \, dx \bigg| \leq \sum_{y \in \mathcal{Y}} \int\limits_{[y,y_+]}\bigg|\frac{\partial^s}{\partial x_1 \cdots \partial x_s} f(x)\bigg| \, dx
\end{equation*}

Since for a ladder $\mathcal{Y}$ of a hyperrectangle $[a,b]$ the equality
$\bigcup\limits_{y\in\mathcal{Y}}[y,y_+] = [a,b]$ holds, the sum builds
the integral over the whole hyperrectangle $[a,b]$:

\begin{equation*}
    \sum_{y \in \mathcal{Y}} \int\limits_{[y,y_+]}\bigg|\frac{\partial^s}{\partial x_1 \cdots \partial x_s} f(x)\bigg| \, dx = \int\limits_{[a,b]}\bigg|\frac{\partial^s}{\partial x_1 \cdots \partial x_s} f(x)\bigg| \, dx.
\end{equation*}

Therefore

\begin{equation*}
    V_{[a,b]}(f) = \sup_{\mathcal{Y}\in\mathbb{Y}} V_{\mathcal{Y}}(f) \leq \int\limits_{[a,b]} \left| \frac{\partial^s}{\partial x_1 \cdots \partial x_s} f(x) \right| \, dx.
\end{equation*}
\end{proof}

\begin{remark} \ \\
    As computing the Hardy--Krause variation is not tractable, the
    following recursively defined upper bound follows immediately from the
    preceding results by Corollary~\ref{cor:variation-approximation}:

    \begin{equation}
        \label{eq:variation_hardy_krause_bound}
        V_{\mathrm{HK}}(f) \leq \widehat{V}_{\mathrm{HK}}(f) := \int\limits_{[a,b]} 
        \left| \frac{\partial^s}{\partial x_1 \cdots \partial x_s} f(x) \right| \, 
        dx + \sum\limits_{i=1}^s \widehat{V}_{\mathrm{HK}}(f|_{x_i=b_i}).
    \end{equation}

    Here, $f|_{x_i=b_i}$ denotes the restriction of $f$ to the boundary with
    $x_i=b_i$. In the first place this is a restriction to the $s-1$ dimensional
    case. In the last recursive step, the univariate case is reached with

    \begin{equation*}
        \widehat{V}_{\mathrm{HK}}(f\big|_{x_i=b_i, \forall i \neq j}) = \int\limits_{[a_j,b_j]} 
        \left| \frac{\partial f\big|_{x_i=b_i, \forall i\neq j}(x_j)}{\partial x_j} \right| \, dx.
    \end{equation*}
    Hereby,  the restricted function $f\big|_{x_i=b_i, \forall i \neq j}$ is
    limited to the one remaining dimension $j$.
\end{remark}


\section{Error estimation of DNN surrogates obtained by \\ low-discrepancy training data}

Given the computable training error $\mathcal{E}_T$ depending on the training
data, we now seek to estimate the generalization error $\mathcal{E}_G$. To get there, we will use an estimate of the generalization gap as in \cite{mishra2021enhancing}.

\begin{theorem} \ \\
    For the surrogate models $\hat{f}$ obtained by \ac{dnn} training of a
    mapping $f$ on the hyperrectangle $[0,1]^s$ with a low-discrepancy sequence
    $\mathcal{S}=\{x_n\}_{n=1}^N\subset [0,1]^s$ yields a generalization gap
    estimated by
    \begin{equation*}
        |\mathcal{E}_G - \mathcal{E}_T| \leq C \cdot \frac{V_{\mathrm{HK}}(|f-\hat{f}|)(\log N)^s}{N},
    \end{equation*}
    where $C$ is a constant and $V_{\mathrm{HK}}(|f-\hat{f}|)$ is the
    Hardy--Krause variation from Definition~\ref{def:variation_hardy_krause}.
\end{theorem}

\begin{proof} \ \\
    This simply follows from the Koksma-Hlawka inequality
    (Theorem~\ref{thm:koksma-hlawka}) applied to the difference $f-\hat{f}$:
    \begin{align*}
        |\mathcal{E}_G - \mathcal{E}_T| &= \bigg| \int\limits_{[0,1]^s} f(x) - \hat{f}(x) \, dx - \frac{1}{N} \sum_{n=1}^N f(x_n) - \hat{f}(x_n) \bigg| \\
        &\leq V_{\mathrm{HK}}(|f-\hat{f}|) \cdot D^*_N \\
        &\leq C \cdot \frac{V_{\mathrm{HK}}(|f-\hat{f}|)(\log N)^s}{N}
    \end{align*}
    Here, $D^*_N$ is the star-discrepancy of the low-discrepancy sequence
    $\mathcal{S}$ and $C$ is a constant depending on the discrepancy of the
    sequence.
\end{proof}

\begin{theorem} \ \\
\label{thm:main-theorem}
Let $f : [0,1]^s \to \mathbb{R}$ with $f\in C^s$ be a map with finite
Hardy--Krause variation, i.e. $\widehat{V}_{HK}(f) < \infty$. Further let
$\hat{f}$ be the \ac{dnn} surrogate with an activation function $\sigma \in C^s$
trained on a low-discrepancy sequence $\mathcal{S}=\{x_n\}_{n=1}^N\subset
[0,1]^s$. Then for a given tolerance $\delta > 0$ the following holds:
\begin{equation*}
    \mathcal{E}_G \leq \mathcal{E}_T + C \cdot \frac{(\log N)^s}{N} + 2\delta.
\end{equation*}
The constant $C$ depends on the learing algorithm and the tolerance $\delta$
associated with the mollifier introduced in the proof.
\end{theorem}

To proof this theorem, we will use the following proposition stating a
multivariate version of the Faà di Bruno formula from
\cite{hardy2006combinatorics}.

\label{prop:multivariate_faa_di_bruno}
\begin{prop}[Multivariate Faà di Bruno formula] \ \\
Let $h \in C^s(\mathbb{R}^s, \mathbb{R}^{d_1})$ and $g \in C^s(\mathbb{R}^{d_1},\mathbb{R}^{d_2})$ be two
functions. Then the following derivative of the composition $g \circ
h$ follows:
\begin{equation*}
    \frac{\partial^s}{\partial x_1 \cdots \partial x_s} (g \circ h)(x) = \sum_{\pi} g^{(|\pi|)}\circ h(x) \cdot \prod_{B\in\pi} \frac{\partial^{|B|}}{\prod_{j\in B}\partial x_j}h(x)
\end{equation*}
Hereby the sum is taken over all partitions $\pi$ of the set $\{1,\ldots,s\}$
and $B$ are the sets within the partition $\pi$.
\end{prop}
A proof of the Multivariate Faà di Bruno formula can be found in
\cite[Proposition 1]{hardy2006combinatorics}.

The multivariate Faà di Bruno formula will be used to proof the following result on the conditions for boundedness of the Hardy--Krause variation of compositions of functions.

% \begin{remark}\label{rem:compositions} \ \\
% To check in what cases the variation in the sense of Vitali is bounded for the
% composition of two functions, we first apply
% Corollary~\ref{cor:variation-approximation} on the composition of two functions
% $h \in C^s([0,1]^s, \mathcal{X})$ and $g \in C^s(\mathcal{X})$, where
% $\mathcal{X}$ is compact. For $g \circ h$ this yields
% \begin{equation*}
%     V_{[0,1]^s}(g \circ h) \leq \int\limits_{[0,1]^s} \bigg| \frac{\partial^s}{\partial x_1 \cdots \partial x_s} (g \circ h)(x) \bigg| \, dx.
% \end{equation*}
% By applying the multivariate Faà di Bruno formula from
% Proposition~\ref{prop:multivariate_faa_di_bruno} it follows that
% \begin{equation*}
%     V_{[0,1]^s}(g \circ h) \leq \int\limits_{[0,1]^s} \bigg| \sum_{\pi} g^{(|\pi|)}\circ h(x) \cdot \prod_{B\in\pi} \frac{\partial^{|B|}}{\prod_{j\in B}\partial x_j}h(x) \bigg| \, dx.
% \end{equation*}
% Now we want to use this inequality to find conditions, such that the variation
% in the sense of Vitali is bounded. For finite dimensions $s$, the number of
% partitions $\pi$ is finite and therefore the sum over all partitions $\pi$ of
% $\{1,\dots,s\}$ is finite. Thus, it suffices to show that each summand is
% bounded.
% \begin{equation*}
%     \int\limits_{[0,1]^s} \bigg| g^{(|\pi|)}\circ h(x) \cdot \prod_{B\in\pi} \frac{\partial^{|B|}}{\prod_{j\in B}\partial x_j}h(x) \bigg| \, dx < \infty.
% \end{equation*}
% As $g \in C^s(\mathcal{X})$, $g^{(\pi)}$ is continuous. As $h$ is also continuous and the integral is over a compact set, the image $g^{(|\pi|)}\circ h([0,1]^s)$ is also compact and therefore bounded. Thus, it suffices to show the following for all partitions $\pi$ of the set $\{1,\ldots,s\}$:
% \begin{equation*}
%     \int\limits_{[0,1]^s} \prod_{B\in\pi} \bigg| \frac{\partial^{|B|}}{\prod_{j\in B}\partial x_j}h(x) \bigg| \, dx < \infty.
% \end{equation*}
% \end{remark}

% As the Variation in the sense of Hardy--Krause is defined by a finite sum over
% variations in the sense of Vitali, the following Corollary follows immediately
% for finite dimensions $s$.

% \begin{cor} \ \\
% \label{cor:bounded_composition}
%     Let $f \in C^s([0,1]^s, \mathcal{X})$ for $\mathcal{X}\subset \mathbb{R}^d$
%     and $s$ finite. If
%     \begin{equation*}
%         \int\limits_{[0,1]^{|v|}} \prod_{B\in\pi} \bigg| \frac{\partial^{|B|}}{\prod_{j\in B} \partial x_j}f(x) \bigg| \, dx < \infty
%     \end{equation*}
%     holds for all partitions $\pi$ of the set $\{1,\ldots,s\}$, then the
%     Hardy--Krause variation of the composition $\eta \circ f$ is bounded for all
%     $\eta \in C^s(\mathcal{X})$:
%     \begin{equation*}
%         V_{\mathrm{HK}}(\eta \circ f) < \infty.
%     \end{equation*}
% \end{cor}

\begin{cor} \ \\
\label{cor:bounded_composition}
Let $h: [a,b]\to\mathcal{X}$ for a hyperrectangle $[a,b]\in\mathbb{R}^s$ and a
compact subset $\mathcal{X}\subset \mathbb{R}$. Further assume that
$\widehat{V}_{\mathrm{HK}}(h) < \infty$. Then, the Hardy--Krause variation of
the composition $g \circ h$ is bounded
\begin{equation*}
    V_{\mathrm{HK}}(g \circ h) < \infty
\end{equation*}
for all $g \in C^s(\mathbb{R})$.
\end{cor}

\begin{proof} \ \\
To check the boundedness of the Hardy--Krause variation of a composition of two
functions, its definition \eqref{eq:hardy-krause-variation} is applied.
\begin{equation*}
    V_{\mathrm{HK}}(g \circ h) = \sum_{v \subset \{1,\ldots,s\}} V_{[a^{-v},b^{-v}]}(g \circ h) (x^{-v};b^{v})
\end{equation*}

As the dimension $s$ is finite, the sum in the definition of the Hardy--Krause
variation has finitely many summands and therefore it suffices to show that each
summand is bounded. It remains to proof that for each subset $v \subset
\{1,\ldots,s\}$ the variation in the sense of Vitali is bounded.

By applying Corollary~\ref{cor:variation-approximation} on the summand, it would
suffice to show  that

\begin{equation*}
    \int\limits_{[a^{-v},b^{-v}]} \bigg| \frac{\partial^{|v|}}{\prod_{j\in v}\partial x_j} (g \circ h)(x^{-v};b^{v}) \bigg| \, dx^{-v} <\infty.
\end{equation*}

Applying the multivariate Faà di Bruno formula
(Proposition~\ref{prop:multivariate_faa_di_bruno}) yields

\begin{equation*}
    \int\limits_{[a^{-v},b^{-v}]} \bigg| \sum_{\pi} \big(g^{(|\pi|)}\circ h\big)(x^{-v};b^{v}) \cdot \prod_{B\in\pi} \frac{\partial^{|B|}}{\prod_{j\in B}\partial x_j}h(x^{-v};b^{v}) \bigg| \, dx^{-v} < \infty
\end{equation*}

with the sum taken over all partitions $\pi$ of the subset $v$. The boundedness
of $h$ which follows from the compactness of $\mathcal{X}$ and the continuity of
the derivative  $g^{(|\pi|)}$ yields an upper bound on $\big(g^{(|\pi|)}\circ
h\big)(x^{-v};b^{v})$.

As the number of partitions $\pi$ of the finite subset $v$ of $\{1,\ldots,s\}$
it remains to show:
\begin{equation*}
    \int\limits_{[a^{-v},b^{-v}]} \prod_{B\in\pi} \bigg| \frac{\partial^{|B|}}{\prod_{j\in B}\partial x_j}h(x^{-v};b^{v}) \bigg| \, dx^{-v} < \infty 
\end{equation*}
for all partitions $\pi$ of the subset $v$. This follows by the definition
\eqref{eq:variation_hardy_krause_bound} of the assumption that
$\widehat{V}_{\mathrm{HK}}(h) < \infty$.
\end{proof}

Now we can prove the theorem.
\begin{proof}[Proof (Theorem~\ref{thm:main-theorem})] \ \\
    First, we define a mollifer function $\eta_\delta \in C^s(\mathbb{R},
    \mathbb{R}_+)$ fullfilling
    \begin{equation*}
        \max \big\{\big\| |a| - \eta_\delta(a) \big\|_\infty\big\} \leq \delta.
    \end{equation*}

    By further defining approximating the the generalization error
    $\mathcal{E}_G$ and the training error $\mathcal{E}_T$ by using the
    differentiable mollifier $\eta_\delta$ instead of the norm
    $\|\cdot\|_\infty$ we obtain
    \begin{align*}
        \mathcal{E}_G^\delta &:= \int_X \eta_\delta(f(x) - \hat{f}(x)) \\
        \mathcal{E}_T^\delta &:= \frac{1}{N} \sum_{n=1}^N \eta_\delta(f(x_n) - \hat{f}(x_n)).
    \end{align*}

    Next to show, is that the mollified error $\eta_\delta\circ(f-\hat{f})$ has
    bounded Hardy--Krause variation.
    
    Since $\hat{f}$ is a composition of affine and activation functions $\sigma
    \in C^s$ by definition \eqref{eq:dnn-definition}, $\hat{f} \in C^s([0,1]^s,
    \mathbb{R})$ holds. Let $\mathbb{1}_s$ denote the identity function on
    $\mathbb{R}^s$. $\widehat{V}_\mathrm{HK}(\mathbb{1}_s) < \infty$ holds
    trivially. Thus applying Corollary~\ref{cor:bounded_composition} for $f = f
    \circ \mathbb{1}_s$ yields

    \begin{equation*}
        {V}_{\mathrm{HK}}(\hat{f}) \equiv \widehat{V}_{\mathrm{HK}}(\hat{f}) <
        \infty.
    \end{equation*}

    By the definition \eqref{eq:variation_hardy_krause_bound} and from the given
    assumption $\widehat{V}_{\mathrm{HK}}(f) < \infty$.
    
    \begin{equation*}
        \widehat{V}_{\mathrm{HK}}(f-\hat{f}) \leq \infty
    \end{equation*}

    holds. Applying Corollary~\ref{cor:bounded_composition} again on $\eta_\delta \circ (f-\hat{f})$ yields 
    
    \begin{equation*}
    V_{\mathrm{HK}}\big(\eta_\delta\circ(f-\hat{f})\big) \leq \widehat{V}_{\mathrm{HK}}\big(\eta_\delta\circ(f-\hat{f})\big) \leq C_1 < \infty.
    \end{equation*}
    
    % Therefore, we apply Corollary~\ref{cor:bounded_composition} and need to
    % therefore to show, that $f-\hat{f} \in C^s([0,1]^s, \mathbb{R})$.

    % The ground truth function $f$ fullfils this by assumptions. As the
    % composition $\hat{f}$ from \eqref{eq:dnn-definition} is a composition of
    % affine and activation functions fullfilling $C^s$ by assumption, it follows
    % that $\hat{f} \in C^s([0,1]^s, \mathbb{R})$.

    % Furthermore, to apply Corollary~\ref{cor:bounded_composition} we need to show that

    % \begin{equation} \label{eq:variation_requirement}
    %     \int\limits_{[0,1]^s} \big| \prod_{B\in\pi} \frac{\partial^{|B|}}{\prod_{j\in B}\partial x_j} (f-\hat{f})(x) \big| \, dx < \infty.
    % \end{equation}

    % holds for all partitions $\pi$ of the set $\{1,\ldots,s\}$.

    % As $f$ has bounded Hardy--Krause variation by assumption and the surrogate
    % model $\hat{f}$ is a composition of affine (hence $C^s$) and sufficiently
    % smooth activation functions $\sigma$, \eqref{eq:variation_requirement} holds
    % for all partitions $\pi$ of the set $\{1,\ldots,s\}$:
    % \begin{equation}
    %     \label{eq:bounded_mollified_generalization_gap}
    %     \eta_\delta(f(x) - \hat{f}(x)) \leq C < \infty \qquad \text{for all } x \in [0,1]^s.
    % \end{equation}

    % The constant $C$ depends on both the dimension $s$ and the tolerance
    % $\delta$ associated with the mollifier. Since it has been established that
    % the mollified difference $\eta_\delta(f(x) - \hat{f}(x))$ is bounded,
    % Remark~\ref{rem:compositions} implies that its Hardy--Krause variation is
    % also bounded.


    for a constanct $C_1$. Therefore, by applying the Koksma-Hlawka inequality
    (Theorem~\ref{thm:koksma-hlawka}) one gets
    \begin{equation}
        \label{eq:hk_generalization_gap}
        |\mathcal{E}_G^\delta - \mathcal{E}_T^\delta| \leq C_2 \cdot \frac{V_{\mathrm{HK}}(\eta_\delta(f-\hat{f}))(\log N)^s}{N}.
    \end{equation}

    With $C=C_1 \cdot C_2$, where the constant $C_1$ depends on $\delta$ and the dimension $s$, it follows that

    \begin{equation}
        |\mathcal{E}_G^\delta - \mathcal{E}_T^\delta | \leq C \cdot \frac{(\log N)^s}{N}.
    \end{equation}

    By construction of $\eta_\delta$ it holds that
    \begin{equation}
        \big| \mathcal{E}_G - \mathcal{E}_G^\delta \big|\leq \delta \qquad \text{and} \qquad \big| \mathcal{E}_T - \mathcal{E}_T^\delta \big| \leq \delta.
    \end{equation}
    Applying by applying triangle inequality to the initial generalization gap, it follows that
    \begin{align*}
        |\mathcal{E}_G - \mathcal{E}_T| &\leq \underbrace{|\mathcal{E}_G - \mathcal{E}_G^\delta|}_{\leq \delta} + \underbrace{|\mathcal{E}_T^\delta - \mathcal{E}_T|}_{\leq \delta} + |\mathcal{E}_G^\delta - \mathcal{E}_T^\delta| \\
        &\leq 2\delta + |\mathcal{E}_G^\delta - \mathcal{E}_T^\delta| \\
        &\leq 2\delta + C \cdot \frac{(\log N)^s}{N}.
    \end{align*}
    It follows the bound on the generalization error as stated in the theorem:
    \begin{equation}
        \mathcal{E}_G \leq \mathcal{E}_T + C \cdot \frac{(\log N)^s}{N} + 2\delta.
    \end{equation}

\end{proof}




















